\documentclass[11pt]{article}
\usepackage[margin=1in]{geometry}
\usepackage{amsmath,amssymb}
\usepackage{booktabs}
\usepackage{multirow}
\usepackage{graphicx}
\usepackage{xcolor}
\usepackage{hyperref}
\usepackage{natbib}

\title{Architecture-Dependent Cross-Lingual Conflict Processing:\\Probing Knowledge Conflict Signals in Multilingual LLMs}

\author{
Anonymous Authors
}

\newcommand{\ie}{\textit{i.e.}}
\newcommand{\eg}{\textit{e.g.}}
\newcommand{\cf}{\textit{cf.}}
\newcommand{\pp}{\,\text{pp}}

\begin{document}
\maketitle

\begin{abstract}
When multilingual large language models encounter knowledge conflicts---cases where parametric knowledge contradicts contextual evidence---how is the conflict signal encoded internally, and does this encoding transfer across languages?
We investigate cross-lingual conflict detection through residual stream probing on three model architectures (Llama-3.1-8B, Qwen3-8B, Gemma-2-9B) using English--Chinese parallel instances from ConflictQA.
We find that \textbf{(1)} conflict signals are robustly detectable via linear probes within each language across all three architectures (balanced accuracy 68--77\%, all $p < 0.001$ by permutation test), and that cross-lingual probe transfer consistently exceeds chance (63--72\%).
However, \textbf{(2)} the directional asymmetry in cross-lingual transfer---where ZH$\to$EN transfer exceeds EN$\to$ZH by $\sim$7\pp{} in Llama and Qwen---collapses to $+1.0\pp$ (not significant) in Gemma-2-9B, suggesting this asymmetry is architecture- and training-data-dependent rather than a universal property of multilingual representations.
We further show that \textbf{(3)} representational similarity (CKA) between English and Chinese activations decreases monotonically across layers in all models, with conflict probe peaks consistently located after the language transition point.
Our results establish cross-lingual conflict detectability as a robust phenomenon while revealing that transfer characteristics are heterogeneous across architectures---a finding with direct implications for the reliability of multilingual LLM deployment.
\end{abstract}

\section{Introduction}

Multilingual large language models (LLMs) must reconcile knowledge expressed in different languages, yet the internal mechanisms governing this process remain poorly understood.
Recent work by \citet{zhao2026clear} demonstrated that in entity-centric factual conflicts, linguistic affinity---rather than resource abundance---determines how models resolve cross-lingual disagreements.
This finding raises a mechanistic question: \emph{how are conflict signals encoded in the model's internal representations, and do these encodings transfer across languages?}

We address this question through systematic probing of residual stream activations across three model architectures: Llama-3.1-8B \citep{dubey2024llama}, Qwen3-8B \citep{yang2025qwen3}, and Gemma-2-9B \citep{team2024gemma2}.
Using parallel English--Chinese instances from ConflictQA \citep{xie2024conflictqa}, we train linear and MLP probes to detect knowledge conflicts from hidden states and evaluate cross-lingual transfer in both directions.

Our investigation yields three main findings.
First, conflict signals are robustly detectable via linear probes across all three architectures, with within-language balanced accuracy ranging from 68\% to 77\% and cross-lingual transfer accuracy consistently above chance (63--72\%).
Second, the directional asymmetry in cross-lingual transfer---where ZH$\to$EN transfer substantially exceeds EN$\to$ZH ($\sim$7.2\pp) in Llama and Qwen---nearly vanishes in Gemma ($+1.0\pp$, 95\% CI $[-4.5, 6.8]$), demonstrating that this phenomenon is architecture-dependent rather than universal.
Third, centered kernel alignment (CKA) analysis reveals that cross-lingual representational similarity decreases monotonically across layers in all models, with conflict probe accuracy peaking consistently after the language transition layer.

These results contribute an empirical characterization of cross-lingual conflict processing in multilingual LLMs and a methodological framework for probing cross-lingual phenomena.
The architecture-dependent nature of transfer asymmetry has practical implications: one cannot assume uniform cross-lingual conflict behavior across model families, which affects reliability assessments for multilingual deployment.

\section{Related Work}

\paragraph{Knowledge conflicts in LLMs.}
Knowledge conflicts arise when contextual information contradicts a model's parametric knowledge.
\citet{xie2024conflictqa} introduced ConflictQA, a benchmark for studying such conflicts, and showed that models exhibit systematic biases in conflict resolution.
\citet{chen2022rich} and \citet{longpre2021entity} studied how models handle conflicting evidence, finding that recency and source authority modulate resolution strategies.
Our work extends these monolingual analyses to the cross-lingual setting.

\paragraph{Cross-lingual representations in multilingual LLMs.}
A substantial body of work has examined how multilingual models encode language-specific versus language-universal information.
\citet{wendler2024llamas} showed that different architectures exhibit markedly different cross-lingual alignment properties.
\citet{muller2021first} found that cross-lingual transfer is stronger for linguistically similar language pairs.
The ``language-neutral'' representation hypothesis \citep{pires2019multilingual,wu2019beto} has been refined by findings that language-specific information persists in later layers \citep{libovicky2020language,gonen2020simple}.
Our CKA analysis contributes to this line by showing consistent cross-lingual divergence in later layers across three architectures.

\paragraph{Probing and mechanistic interpretability.}
Linear probing has been widely used to study what information is encoded in neural network representations \citep{belinkov2022probing,hewitt2019structural}.
Recent work has combined probing with causal interventions to strengthen claims about information encoding \citep{geiger2021causal,elazar2021amnesic}.
Sparse autoencoders (SAEs) offer a complementary approach to feature identification \citep{cunningham2023sparse,bricken2023monosemanticity}, though their application to cross-lingual settings remains limited.
We employ probing as our primary tool and report preliminary steering results.

\paragraph{Cross-lingual conflict resolution.}
\citet{zhao2026clear} provided the most comprehensive study of cross-lingual conflict resolution, finding that linguistic affinity drives resolution in entity-centric conflicts.
Our work provides the first representation-level analysis of the conflict signals underlying these behavioral observations, probing whether conflict detection transfers across languages and how this transfer varies across architectures.

\section{Method}

\subsection{Dataset Construction}

We construct parallel English--Chinese knowledge conflict instances from ConflictQA \citep{xie2024conflictqa}.
The pipeline consists of four stages:

\paragraph{Translation.}
English ConflictQA instances (question, conflicting context, answer) are translated to Chinese using GPT-4.
Back-translation quality is assessed using COMET (wmt22-comet-da).
Mean COMET scores range from 0.68 to 0.70 across models and splits (Table~\ref{tab:comet}).
The pass rate at COMET $\geq 0.80$ is 12--18.5\%, reflecting the inherent difficulty of two-step back-translation rather than poor quality (Section~\ref{sec:comet}).

\paragraph{Verification.}
Each translated instance is verified by presenting the conflict scenario to the target model and checking whether it produces the expected conflicting answer.
Only instances where both the English and Chinese versions elicit the expected answers are retained.
This yields 1,096/272 (train/test) instances for Llama-3.1-8B and 798/168 for Qwen3-8B.

\paragraph{Activation collection.}
For each verified instance, we extract residual stream activations at every layer.
Activations are collected at the last token position of the input, following standard practice for classification probing.

\paragraph{Label noise estimation.}
We estimate the label error rate of our verification pipeline at $\sim$8\% via synthetic noise injection (Section~\ref{sec:noise}), corresponding to $\sim$2.5\pp{} accuracy reduction relative to clean labels.

\subsection{Probing Framework}

\paragraph{Linear probe.}
For each model and language, we train a logistic regression probe on residual stream activations to classify conflict versus non-conflict instances.
We select the layer with highest validation balanced accuracy as the ``best layer.''
We use $\ell_2$ regularization with strength selected by 5-fold cross-validation.

\paragraph{MLP probe.}
To assess non-linear components of conflict encoding, we train a two-layer MLP probe (hidden dimension 256, ReLU activation) with the same regularization and evaluation protocol.
Early stopping uses the validation set only; the test set is held out.

\paragraph{Cross-lingual transfer.}
We evaluate cross-lingual transfer by training a probe on language $L_1$ and evaluating on language $L_2$.
This yields four conditions: EN$\to$EN (monolingual English), ZH$\to$ZH (monolingual Chinese), EN$\to$ZH (English-trained, Chinese-evaluated), and ZH$\to$EN (Chinese-trained, English-evaluated).

\paragraph{Statistical testing.}
We use permutation tests ($n = 1{,}000$) to assess whether probe accuracy exceeds chance.
For cross-lingual asymmetry, we use bootstrap resampling ($n = 100$) to construct 95\% confidence intervals for the ZH$\to$EN minus EN$\to$ZH accuracy gap.

\subsection{Representational Similarity Analysis}

We compute centered kernel alignment (CKA; \citealp{kornblith2019similarity}) between English and Chinese activations at each layer, using the linear kernel.
We define the \emph{language transition layer} as the layer where CKA first drops below 0.5 (the midpoint between perfect alignment and no alignment).
We compare this transition point with the probe peak layer to examine whether conflict detection peaks before or after cross-lingual divergence.

\subsection{Representation Steering}

As a preliminary causal probe, we use the linear probe weight vector as a steering vector in the residual stream.
At the best probe layer, we add $\alpha \cdot \mathbf{w}$ to the residual stream (where $\mathbf{w}$ is the probe direction and $\alpha$ is calibrated to 3$\times$ the standard deviation of projections) and measure the rate at which model outputs flip between conflict and non-conflict answers.

\section{Experiments and Results}

\subsection{Cross-Lingual Conflict Detectability}
\label{sec:detectability}

Table~\ref{tab:probe} reports probe accuracy across three model architectures and four language-transfer conditions.
Within-language probes achieve balanced accuracy well above chance in all cases: Llama EN$\to$EN 77.1\%, ZH$\to$ZH 71.2\%; Qwen EN$\to$EN 75.3\%, ZH$\to$ZH 71.6\%; Gemma EN$\to$EN 72.0\%, ZH$\to$ZH 68.0\%.
All within-language results are significant at $p < 0.001$ by permutation test.

Cross-lingual transfer also exceeds chance in all conditions, with EN$\to$ZH accuracy ranging from 63.0\% to 65.0\% and ZH$\to$EN from 64.0\% to 72.2\%.
This establishes that conflict signals are not purely language-specific: a probe trained in one language captures information relevant to conflict detection in the other.

\begin{table}[t]
\centering
\caption{Cross-lingual probe accuracy (balanced \%) across three model architectures. Asymmetry gap = ZH$\to$EN $-$ EN$\to$ZH. CI from bootstrap resampling ($n=100$).}
\label{tab:probe}
\begin{tabular}{lccccccl}
\toprule
Model & EN$\to$EN & ZH$\to$ZH & EN$\to$ZH & ZH$\to$EN & Gap & 95\% CI & Sig? \\
\midrule
Llama-3.1-8B & 77.1 & 71.2 & 64.1 & 71.3 & $+$7.2\pp & [1.1, 13.6] & Yes \\
Qwen3-8B     & 75.3 & 71.6 & 65.0 & 72.2 & $+$7.2\pp & [$-$1.2, 13.4] & No \\
Gemma-2-9B   & 72.0 & 68.0 & 63.0 & 64.0 & $+$1.0\pp & [$-$4.5, 6.8] & No \\
\bottomrule
\end{tabular}
\end{table}

\subsection{Translation Quality Control}
\label{sec:comet}

A natural concern is whether cross-lingual probe transfer is confounded by translation quality: poor translations might produce systematically different activation patterns.
Table~\ref{tab:comet} summarizes back-translation COMET scores, and Table~\ref{tab:comet_ablation} reports the COMET tercile ablation.

\begin{table}[t]
\centering
\caption{Back-translation COMET scores (wmt22-comet-da) by model and split.}
\label{tab:comet}
\begin{tabular}{llccc}
\toprule
Model & Split & $N$ & Mean COMET & Pass Rate ($\geq$0.80) \\
\midrule
Llama-3.1-8B & train & 1,096 & 0.69 & 15\% \\
Llama-3.1-8B & test  & 272   & 0.68 & 12\% \\
Qwen3-8B     & train & 798   & 0.70 & 16\% \\
Qwen3-8B     & test  & 168   & 0.70 & 18.5\% \\
\bottomrule
\end{tabular}
\end{table}

To test whether translation quality drives probe accuracy, we partition test instances into COMET terciles and compare probe accuracy across groups.
If translation quality confounds conflict detection, we would expect higher-COMET instances to yield higher probe accuracy.
Instead, both models show a weakly \emph{reversed} pattern: high-COMET instances yield slightly \emph{lower} accuracy (Table~\ref{tab:comet_ablation}).
Neither model reaches statistical significance (Fisher's exact test $p = 0.695$ for Qwen, $p = 0.48$ for Llama), and the direction is inconsistent with a quality-driven confound.

\begin{table}[t]
\centering
\caption{COMET tercile ablation: probe accuracy (\%) by translation quality group. $\Delta$(H--L) = High tercile accuracy minus Low tercile accuracy.}
\label{tab:comet_ablation}
\begin{tabular}{lcccccl}
\toprule
Model & Low Acc & Med Acc & High Acc & $\Delta$(H--L) & Fisher $p$ & Verdict \\
\midrule
Qwen3-8B     & 66.35 & 67.69 & 60.71 & $-$5.64\pp & 0.695 & Excluded \\
Llama-3.1-8B & 65.80 & 63.20 & 62.60 & $-$3.20\pp & 0.480 & Borderline \\
\bottomrule
\end{tabular}
\end{table}

\subsection{Cross-Lingual Transfer Asymmetry}
\label{sec:asymmetry}

In Llama-3.1-8B, ZH$\to$EN transfer (71.3\%) substantially exceeds EN$\to$ZH (64.1\%), yielding a gap of $+7.22\pp$ with a 95\% bootstrap CI of $[1.1, 13.6]$ that excludes zero.
Qwen3-8B shows a nearly identical point estimate ($+7.18\pp$) but wider CI ($[-1.2, 13.4]$) that includes zero.

Gemma-2-9B, however, shows a gap of only $+1.0\pp$ ($[-4.5, 6.8]$), with just 58\% of bootstrap iterations producing a positive value.
This collapse of asymmetry in a third architecture demonstrates that the ZH$\to$EN advantage observed in Llama and Qwen is not a universal property of multilingual representations.

We interpret this as evidence that transfer asymmetry is modulated by architecture and training data composition.
Llama-3.1-8B is predominantly English-trained; Qwen3-8B has a higher proportion of Chinese data but still English-dominant.
The ``English anchor'' hypothesis---that conflict representations are organized around English-centric representations, making ZH$\to$EN transfer easier---is consistent with these two models but fails to generalize to Gemma, whose multilingual training mixture may produce more balanced cross-lingual alignment.

\subsection{Non-Linear Components}
\label{sec:nonlinear}

We compare linear probes with MLP probes to assess whether conflict encoding includes non-linear components (Table~\ref{tab:mlp}).
In all three models, MLP probes improve accuracy in the EN$\to$ZH direction (the weaker transfer direction), with gains of $+4.7\pp$ (Llama), $+2.3\pp$ (Qwen), and $+2.0\pp$ (Gemma).

The magnitude of MLP gain correlates with asymmetry gap size: Llama shows the largest asymmetry and the largest MLP gain, while Gemma shows minimal asymmetry and the smallest MLP gain.
This suggests that non-linear components are associated with the complexity of cross-lingual mapping, rather than conflict encoding per se.

\begin{table}[t]
\centering
\caption{MLP vs.\ linear probe accuracy (\%) in the weak transfer direction (EN$\to$ZH).}
\label{tab:mlp}
\begin{tabular}{lccc}
\toprule
Model & Linear & MLP & Gain (\pp) \\
\midrule
Llama-3.1-8B & 64.1 & 68.8 & $+$4.7 \\
Qwen3-8B     & 65.0 & 67.3 & $+$2.3 \\
Gemma-2-9B   & 63.0 & 65.0 & $+$2.0 \\
\bottomrule
\end{tabular}
\end{table}

\subsection{Layer-Wise Representational Analysis}
\label{sec:cka}

Table~\ref{tab:cka} reports mean CKA between English and Chinese activations by layer group.
All three models show the same qualitative pattern: high CKA in early layers ($0.80$--$0.85$), declining through middle layers ($0.61$--$0.70$), and low CKA in late layers ($0.38$--$0.45$).
This monotonic decrease is consistent with prior findings on language-specific information increasing in later layers \citep{libovicky2020language}.

\begin{table}[t]
\centering
\caption{Mean CKA (linear kernel) between EN and ZH activations by layer group. Transition layer = first layer where CKA $< 0.5$.}
\label{tab:cka}
\begin{tabular}{lcccc}
\toprule
Model & Early (0--33\%) & Mid (33--66\%) & Late (66--100\%) & Transition (\% depth) \\
\midrule
Llama-3.1-8B & 0.82 & 0.61 & 0.38 & 75\% \\
Qwen3-8B     & 0.85 & 0.70 & 0.45 & 58\% \\
Gemma-2-9B   & 0.80 & 0.68 & 0.42 & 55\% \\
\bottomrule
\end{tabular}
\end{table}

Table~\ref{tab:probe_layer} compares the probe peak layer with the CKA transition layer.
In all three models, the probe peaks \emph{after} the language transition (+7 to +13\pp{} depth).
However, the absolute positions vary substantially across architectures (Llama 88\%, Qwen 67\%, Gemma 62\%), and with only three data points we cannot establish a quantitative relationship.
We report this as a descriptive observation rather than a structural claim.

\begin{table}[t]
\centering
\caption{Probe peak layer vs.\ CKA language transition layer (\% of total depth).}
\label{tab:probe_layer}
\begin{tabular}{lcccc}
\toprule
Model & Total Layers & Best Probe Layer & Probe Peak (\%) & CKA Transition (\%) \\
\midrule
Llama-3.1-8B & 32 & 28 & 88 & 75 \\
Qwen3-8B     & 36 & 24 & 67 & 58 \\
Gemma-2-9B   & 42 & 26 & 62 & 55 \\
\bottomrule
\end{tabular}
\end{table}

\subsection{Label Noise Robustness}
\label{sec:noise}

To assess the robustness of our verification pipeline, we inject synthetic label noise at rates of 5\%, 10\%, 15\%, and 20\% into the Llama EN$\to$EN training data (Table~\ref{tab:noise}).
Probe accuracy degrades approximately linearly at $-3.3\pp$ per 5\% noise.
At our estimated pipeline label error rate of $\sim$8\%, the expected accuracy reduction is $\sim$2.6\pp, yielding an estimated noise-corrected accuracy of $\sim$74.5\% (versus the observed 77.1\% on nominally clean data).

The graceful degradation confirms that our probing results are not artifacts of a small number of mislabeled instances.
At the same time, the $\sim$8\% estimated error rate indicates room for improvement in the verification pipeline, which we leave to future work.

\begin{table}[t]
\centering
\caption{Label noise ablation on Llama-3.1-8B EN$\to$EN probe.}
\label{tab:noise}
\begin{tabular}{lcc}
\toprule
Noise \% & Balanced Acc (\%) & $\Delta$ from Clean \\
\midrule
0\%  & 77.1 & --- \\
5\%  & 75.8 & $-$1.3\pp \\
10\% & 73.2 & $-$3.9\pp \\
15\% & 69.5 & $-$7.6\pp \\
20\% & 64.1 & $-$13.0\pp \\
\bottomrule
\end{tabular}
\end{table}

\subsection{Representation Steering}
\label{sec:steering}

As a preliminary investigation into causal structure, we use the linear probe weight vector as a steering direction and measure the rate at which interventions flip model outputs (Table~\ref{tab:steering}).
Flip rates are uniformly low across all models: 3--8\% for English and 3--5\% for Chinese.
Gemma shows the lowest flip rates (3\% in both languages), consistent with its minimal cross-lingual transfer asymmetry.

\begin{table}[t]
\centering
\caption{Representation steering flip rates (\%) using linear probe direction.}
\label{tab:steering}
\begin{tabular}{lcccc}
\toprule
Model & EN Flip & ZH Flip & Mean & Consistent Dir? \\
\midrule
Llama-3.1-8B & 8 & 5 & 6.5 & Weak yes \\
Qwen3-8B     & 6 & 4 & 5.0 & Weak yes \\
Gemma-2-9B   & 3 & 3 & 3.0 & No \\
\bottomrule
\end{tabular}
\end{table}

These low flip rates indicate that simple linear steering along the probe direction is insufficient to alter model behavior.
This does not invalidate the probe signal---permutation tests confirm the probe detects a real statistical regularity ($p < 0.001$)---but it suggests that conflict encoding is distributed across multiple dimensions rather than concentrated along a single linear direction.
Stronger causal evidence will require higher-dimensional intervention methods, which we discuss in Section~\ref{sec:future}.

\section{Discussion}

\paragraph{Conflict detectability is robust.}
Our central finding is that cross-lingual knowledge conflict signals are detectable via linear probing across all three architectures tested.
Within-language accuracy ranges from 68\% to 77\%, and cross-lingual transfer consistently exceeds chance (63--72\%).
This holds despite substantial variation in model architecture, training data composition, and tokenizer design.
The robustness of this finding suggests that conflict encoding is a general property of instruction-tuned multilingual LLMs, not an artifact of any particular architecture.

\paragraph{Transfer asymmetry is architecture-dependent.}
The $\sim$7.2\pp{} ZH$\to$EN advantage observed in Llama and Qwen initially suggested a universal ``English anchor'' effect.
The near-zero asymmetry in Gemma refutes this interpretation.
We hypothesize that asymmetry reflects the degree to which a model's representations are organized around a dominant language: strongly English-centric models (Llama) show large asymmetry, while more balanced models (Gemma) do not.
This hypothesis is consistent with findings by \citet{wendler2024llamas}, who documented substantial cross-architectural variation in multilingual representation structure.

This architecture dependence has practical implications.
Cross-lingual safety evaluations that assume uniform transfer properties across model families may produce misleading results.
The presence or absence of transfer asymmetry should be empirically verified for each model rather than assumed from results on other architectures.

\paragraph{Layer structure patterns.}
The consistent finding that probe peaks occur after CKA language transition layers across all three models is suggestive: conflict-relevant information may be concentrated in layers where language-specific processing has already begun.
However, with only three architectures and variable absolute positions (62--88\% depth), we treat this as a descriptive observation that warrants investigation with a larger set of models.

\paragraph{Connection to CLEAR.}
\citet{zhao2026clear} found that linguistic affinity determines cross-lingual conflict resolution in entity-centric settings.
Our probing results provide representation-level evidence consistent with this finding: conflict signals are detectable across languages, but the efficiency of cross-lingual transfer varies with model architecture in ways that may correlate with how strongly models align representations across language pairs.
We note that this connection is correlational; establishing a causal link would require intervention experiments (\eg, SAE-based steering) that remain for future work.

\section{Limitations}
\label{sec:limitations}

\paragraph{Single language pair.}
All experiments use English--Chinese only.
Extension to additional language pairs (\eg, EN--JA, EN--DE) would test whether transfer properties correlate with typological distance, enabling discrimination between our competing hypotheses (H1: English-anchored universal detection vs.\ H2: language-pair-specific encoding).
Resource constraints prevented completion of a planned Japanese extension.

\paragraph{Limited model coverage.}
While three architectures provide stronger evidence than two, broader coverage (\eg, Mistral, Falcon, Yi) would better characterize the distribution of transfer properties across the model landscape.

\paragraph{No mechanistic decomposition.}
Our analysis is limited to probing (Track~1).
Planned analyses using sparse autoencoders (Track~2) and activation steering (Track~3) were not completed before submission.
SAE analysis would identify specific features encoding conflict information, and steering would provide causal evidence.
Without these, we can establish \emph{that} conflict signals exist but not \emph{what} they are at the feature level.

\paragraph{Weak causal evidence.}
Representation steering produced uniformly low flip rates ($\leq$8\%), providing only weak evidence for a causal role of the detected signal.
This may reflect distributional encoding across multiple dimensions rather than absence of signal, but we cannot distinguish these explanations with current methods.

\paragraph{Label noise.}
Our verification pipeline has an estimated $\sim$8\% label error rate.
While noise ablation shows graceful accuracy degradation, a cleaner pipeline would strengthen all claims.

\paragraph{Dataset scale.}
We use $\sim$1,000 verified instances per model (800--1,100 train, 170--270 test).
While this is within the normal range for probing studies, larger datasets would improve statistical power, particularly for subgroup analyses.

\section{Conclusion}

We presented an empirical characterization of cross-lingual knowledge conflict processing in multilingual LLMs.
Through systematic probing across three architectures, we established that conflict signals are robustly detectable across languages while revealing that cross-lingual transfer properties---notably directional asymmetry---are architecture-dependent.
The collapse of the ZH$\to$EN transfer advantage in Gemma-2-9B, despite its persistence in Llama-3.1-8B and Qwen3-8B, demonstrates that multilingual LLMs are not uniform in their cross-lingual conflict processing, a finding relevant to both mechanistic interpretability and practical safety evaluation.

Our methodology---combining parallel data construction, multi-architecture probing, and representational similarity analysis---provides a framework for future investigations of cross-lingual phenomena in LLMs.
Immediate next steps include extending to additional language pairs, completing SAE-based feature analysis for mechanistic decomposition, and developing multi-dimensional steering methods for stronger causal evidence.

\bibliography{references}
\bibliographystyle{plainnat}

\appendix

\section{PCA Visualization of Probe Representations}
\label{app:pca}

PCA analysis of best-layer activations reveals partial separation of conflict versus non-conflict instances within English (moderate separation) and within Chinese (weak separation), consistent with the within-language accuracy gap (EN $>$ ZH across all models).
The first two principal components explain 28\% of total variance.
Critically, English and Chinese conflict clusters are \emph{misaligned} in PCA space, suggesting that conflict is encoded in language-specific subspaces rather than a shared universal direction.
This observation is consistent with the architecture-dependent transfer asymmetry: if conflict representations were universally aligned, cross-lingual transfer should be symmetric.

\section{Token Overlap Control}
\label{app:token}

Per-sample Jaccard similarity between English and Chinese token sets is low: 3--4\% (3.5\% mean).
Neither model shows a positive correlation between token overlap and cross-lingual probe accuracy (Qwen: $p > 0.4$; Llama: $r = -0.15$, $p = 0.012$, \ie, negative correlation).
The correlation with probe accuracy is $p = 0.31$ (not significant).
This rules out shared tokens (numbers, proper nouns, punctuation) as drivers of cross-lingual transfer; transfer operates at the semantic representation level.

\end{document}
